\section{Limitations}
\label{sec:limitations}

\subsection{Internal validity}

The oracle shares the model of the engine: it recomputes the position by the closed form of the stock and finds the root by bisection, so it tests the solver, the bracketing, the rounding and the refusal logic, and not the model. An error in the compounding model itself, for example a wrong treatment of a rate of 1, would be shared by both. The model is checked elsewhere in the product against the printed percentages of the supervisory note and against a day-by-day stock, which this paper does not repeat. The corpus is a plant: the three refusal reasons occur because it is built to produce them, and the rates of its segments are drawn from a distribution chosen by the author, so the shares in \cref{tab:main-results} describe that corpus, in which \kRefused{} of \kInst{} instances are refused, and not the frequency of refusals in practice. The comparators are those a validator would assemble from the forward tools; a more careful scan, with an adaptive grid, would close the gap to the proposed solver in accuracy and not in the handling of refusals. The corpus run took \kWallMin{} minutes, and the experiment was run on one machine, so the timings carry the load of that machine.

\subsection{External validity}

The committed bank is one bank, with real balance-sheet aggregates and synthetic behaviour, valued at one date. The deposit mapping, the monetisation schedule, the calendar and the calibration rule are conventions, and \cref{tab:band} shows that the answer depends on them. The base rates are the transcribed values of one supervisory exercise on a TARGET2 calendar, applied here on a United States business-day calendar. The labelled run is one event with three rows that fall inside the five-day ladder, and the other two 2023 rows have horizons of 90 days, beyond the ladder, so they are not scored. A score on three rows has no sampling interval; the two denominators of the observed outflow are the only uncertainty the paper reports for it. Nothing in the paper concerns what the failed bank did, and the multipliers are not its headroom. The corpus covers constant and two-phase rate paths on horizons of 5 to 125 days; rates that vary through the horizon in other ways, such as a path that accelerates late, are not generated, and the proposition on the exact round trip is not tested on them.

\subsection{Construct validity and scope}

The multiplier is a one-parameter family: it scales every segment by the same factor, so it cannot represent a run that concentrates on one segment. The cap makes the scenario a share but also means that a segment already at a rate of 1, such as the financial segment, does not contribute to the distance. The capacity is exogenous: a faster run does not change haircuts, collateral calls or the willingness of the central bank to lend, all of which a real run would move, and the position is evaluated at the end of the day, so intraday shortfalls are not seen. The target of 30 days is the example of the research that motivates the rung, not a regulatory threshold. The related work on reverse liquidity stress testing contains three book chapters that we could not read, so the novelty of the solver is limited to what this paper defines and measures, and the absence of a published multiplier with a refusal rule is unverified beyond the search made. Finally, the paper does not model run dynamics with a stochastic depositor process, which would give the survival period a distribution and the multiplier a meaning in probability.
